\section{Other errors and the limits of causal attribution}
The explicitly derived models cover resistance mismatch, constant coherent
frame rotation, additive voltage/current/speed perturbations, current relabeling,
and a simple timing extension. Gain errors can be expressed locally as
$\vect\varepsilon_u=\matr G_u\voltage+\vect b_u$ and
$\vect\varepsilon_i=\matr G_i\current+\vect b_i$, but their parity depends
on those matrices and on the frame in which gains and offsets originate.
These expressions parameterize errors; they do not establish a complete
sensor-identification procedure.

Deadtime and semiconductor drops contribute structured terminal-voltage errors
that can depend on phase-current signs, switching conditions, and temperature.
Their dq averages are not generally constant sensor offsets. Work on magnetic
map identification explicitly addresses inverter nonlinearities and uncertain
resistance \cite{liu2018}; here they remain nuisance sources unless an
independent voltage model is available. Commanded inverter voltage and actual
terminal voltage must not be silently equated.

\begin{table}[tb]
\centering\small
\begin{tabularx}{\textwidth}{@{}lXXX@{}}
\toprule
Effect & Forbidden d channel & Forbidden q channel & Speed behavior\\
\midrule
Resistance mismatch & Odd, $-i_q\resistanceerror/\omegae$ & Even, $i_d\resistanceerror/\omegae$;
zero at $i_d=0$ & $1/\omegae$ for fixed mismatch\\
Constant angle & Odd angle sensitivity & Even angle sensitivity & No explicit
speed factor in quasi-static map\\
Saturation alone & None under reflection & None under reflection & Magnetic
map held fixed\\
Saturation with angle & Changes odd sensitivity & Changes even sensitivity &
Through operating point or changing magnetics\\
Paired scalar speed error & None for ideal parity & None for ideal parity &
Common multiplicative flux error\\
Voltage/inverter error & Potentially odd & Potentially even & Voltage error
divided by speed; voltage error itself may vary\\
\bottomrule
\end{tabularx}
\caption{Conditional fingerprints derived earlier. Zero forbidden residual
does not imply zero reconstruction error. Pair speed and thermal state must
match; channel amplitudes may vanish at special operating points.}
\label{tab:fingerprints}
\end{table}

\subsection{Physical alternatives remain open}
An iron-loss branch can make the measured stator current the sum of
magnetizing and loss currents. Reflection of the magnetic law is then a
statement about magnetizing coordinates, and equal/opposite measured-current
keys need not correspond to mirrored magnetic states. This distinction is
explicit in loss-aware measurement and machine models
\cite{richter2014,kruener2019}. Speed-dependent iron-loss currents are one
possible source of residuals in measured-current coordinates. The present
analysis does not identify those currents or prove that iron loss explains
these datasets. Additional torque, power, terminal-voltage and thermal
information would be needed to separate losses from competing effects.
The nonlinear current/speed loss model of Hackl et al. \cite{hackl2021}
provides context rather than an estimator used here.

Dead time, semiconductor drops and reconstructed inverter voltage can change
with current sign, switching conditions and temperature. Offsets and gains
in current and voltage change both reconstructed values and, for current,
the map argument; angle and timing mismatch add the effects derived above.
The exact current-proportional voltage confounder survives any number of
speeds if its coefficient stays constant. A resistance-equivalent residual
therefore need not be an ohmic winding-resistance error even where it is flat.

AC resistance, skin effect, proximity effect, current displacement and
frequency-dependent iron-loss behavior are further possible contributors.
For a homogeneous conductor in the classical local approximation, skin
depth scales as $\sqrt{2/(\omega\mu\sigma)}$: frequency and material
parameters set penetration, while geometry determines its consequence for
current distribution. This formula is not an operating-current law for this
machine. Effective machine AC resistance may additionally depend on field
distribution, saturation and proximity effects. Such dependencies cannot be
assigned to the present residuals without conductor geometry and independent
loss measurements. The basic field and winding-loss background follows
\cite{binder2017}.

Real construction or magnet asymmetry can break reflection physically;
not every forbidden component is a measurement error. Harmonics, skew and
rotor-position averaging require care; symmetric skew need not destroy
an effective reflection. Hysteresis makes flux history-dependent.
Temperature affects winding drops and magnetics. Saturation alone respects
the stated ideal parity, but changes the angle fingerprint through
$\Ldiff$. A constant coherent angle offset may retain reciprocity while
rotating the natural reflection axis. Neither a reciprocity check nor an
absence of forbidden residuals is a complete physical-validity certificate.

\subsection{Experimental and statistical limits}
The measured support is reduced by exact requested-key matching, and the
multi-speed high-current samples contain only six and seven pairs. Their
MAD values describe sampled spread, not general estimator precision.
Requested and measured current coordinates differ; without an independently
validated local magnetic derivative the resulting pairing bias is unbounded
by the present data reduction. References to rotor or stator temperature do
not establish conductor thermal equilibrium. Measured speed is available
but the analysis uses the requested breakpoint, and the provenance of the
averaged voltage signals needs verification. No independent sensor-error
budget or repeated calibrated experiment is supplied here.

Local Delaunay or other mirror interpolation could increase support in a
future study, but would introduce correlated interpolation bias and support
requirements. It is not implemented in this draft. Similarly, derivatives
for angle identification cannot be accepted merely because a flexible model
fits the erroneous measured samples. The present work performs no magnetic
field fit, joint nuisance estimation, formal hypothesis test or unique
physical decomposition. The retained analytical model, real-data residuals
and synthetic sign checks are distinct levels of evidence.

\paragraph{Evidence boundary.}
Demonstrated results are reproducible reconstruction and symmetry residuals
on the supplied files, their spatial structure, comparisons across thermal
references and speed, and analytical/synthetic checks of conditional error
models. The measurements weaken the sole constant-resistance hypothesis as
a global explanation. They do not uniquely distinguish winding resistance,
angle, iron-loss, inverter, sensor, pairing or real-asymmetry contributions.
A local constant equivalent is compatible with a resistance contribution,
but does not prove its presence or quantify its unique share.
